\rhead{DS-EL6: Graph Mining \& Social Network Analysis}
\phantomsection 
\section*{Graph Mining \& Social Network Analysis (DS-EL6)}
\label{elec:DS_6}

\noindent \small{\hyperref[main:scheme]{$\leftarrow$ Back to Scheme}} \vspace{0.5cm}

\noindent \textbf{Credits:} 4 (3-0-2)

\subsection*{Theory}
\noindent\textbf{Unit 1} \\
\textit{Graph Theory Foundations for Analytics:} Graph representations (adjacency matrix, edge list, CSR/CSC), Directed/undirected/multigraphs, Basic metrics (degree, density, diameter, average path length), Connected components and strongly connected components, Graph isomorphism and canonical labeling, Bipartite graphs and projections.

\vspace{0.3cm}

\noindent\textbf{Unit 2} \\
\textit{Network Centrality and Ranking:} Degree centrality variants (in-degree, out-degree, weighted), Closeness centrality and harmonic mean, Betweenness centrality (Brandes algorithm), Eigenvector centrality and Katz matrix, PageRank (power iteration, personalized, topic-sensitive), HITS (authority/hub scores), SALSA improvement.

\vspace{0.3cm}

\noindent\textbf{Unit 3} \\
\textit{Community Detection and Clustering:} Modularity optimization and Louvain method, Spectral clustering (normalized Laplacian), Label propagation algorithm, Infomap (random walks), Clique percolation method, Overlapping communities (CPM, SLPA), Hierarchical clustering (single-linkage dendrograms), Benchmark networks (LFR generator).

\vspace{0.3cm}

\noindent\textbf{Unit 4} \\
\textit{Graph Mining Algorithms:} Frequent subgraph mining (Apriori-based, gSpan), Graph motif discovery and network motifs, Graphlet counting and degree distributions, Triangle counting and enumeration, Graph kernels (graphlet, shortest-path, Weisfeiler-Lehman), Subgraph matching and isomorphism testing, Graph alignment.

\vspace{0.3cm}

\noindent\textbf{Unit 5} \\
\textit{Social Network Analysis and Applications:} Small-world phenomenon and clustering coefficient, Scale-free networks and power-law distributions, Link prediction (common neighbors, Jaccard, Adamic-Adar, Katz, matrix factorization), Influence maximization (greedy algorithm, TIM/CELF), Cascade models (Independent Cascade, Linear Threshold), Signed networks and balance theory, Temporal networks and dynamic analysis.

\newpage

\subsection*{Practical}
\noindent\textbf{Lab 1: Graph Construction and Basic Metrics} \\
Focuses on NetworkX/Graph-tool basics, degree distributions, centrality computation.

\vspace{0.2cm}

\noindent\textbf{Lab 2: PageRank and Ranking Algorithms} \\
Focuses on power iteration implementation, personalized PageRank, convergence analysis.

\vspace{0.2cm}

\noindent\textbf{Lab 3: Community Detection Methods} \\
Focuses on Louvain, spectral clustering, modularity evaluation comparison.

\vspace{0.2cm}

\noindent\textbf{Lab 4: Graph Sampling Techniques} \\
Focuses on random walk, forest fire, node/edge sampling, estimation accuracy.

\vspace{0.2cm}

\noindent\textbf{Lab 5: Frequent Subgraph Mining} \\
Focuses on gSpan, FSG, subgraph enumeration efficiency.

\vspace{0.2cm}

\noindent\textbf{Lab 6: Link Prediction} \\
Focuses on supervised/unsupervised methods, ROC-AUC evaluation, temporal validation.

\vspace{0.2cm}

\noindent\textbf{Lab 7: Influence Maximization} \\
Focuses on greedy algorithm, CELF optimization, simulation-based evaluation.

\vspace{0.2cm}

\noindent\textbf{Lab 8: Graph Neural Networks Basics} \\
Focuses on GCN, GAT, node classification, link prediction tasks.

\vspace{0.2cm}

\noindent\textbf{Lab 9: Dynamic Network Analysis} \\
Focuses on temporal snapshots, streaming algorithms, evolution tracking.

\vspace{0.2cm}

\noindent\textbf{Lab 10: Social Network Analysis Pipeline} \\
Focuses on complete analysis workflow from raw data to insights/visualization.

\vspace{0.2cm}

\newpage
\rhead{Curriculum Schema}
